\section{La base del gerente de producto: perspectiva de tercero, experiencia y datos}

La sección anterior explicó que este episodio debe leerse en clave de juicio de producto; esta sección examina primero de dónde proviene esa capacidad de juicio. En la universidad, Ming Chaoping (明超平) participaba en debates y destacaba sobre todo como cuarto orador; concluyó que lo más importante en un debate no es convencer al adversario, sino conservar siempre la perspectiva de un tercero. Ese entrenamiento se convirtió después en un método de producto: el usuario no es el adversario de las discusiones internas del equipo, sino el jurado que observa desde el público. En el instante en que abre el producto, el usuario no conoce tu historia, tu estrategia ni tu esfuerzo; solo sabe si puede entenderlo y usarlo de inmediato.

\begin{importantbox}{Nota de clase: volverse ingenuo en un segundo}
Ming Chaoping cita la intuición de producto al estilo de Zhang Xiaolong (张小龙): hay que convertirse en un usuario común que entra en el escenario por primera vez. Cualquier software que requiera explicaciones largas, ventanas emergentes y una instrucción compleja para poder usarse no ha completado de verdad su diseño de experiencia.
\end{importantbox}

\subsection{La experiencia no son los datos}

Después de la perspectiva de tercero, esta subsección aborda la primera lección que le dio OnePlus: un producto no es un conjunto de indicadores, sino la sensación continua del usuario en una situación concreta. La pregunta que hay que responder aquí es por qué una función que los datos parecen respaldar puede, aun así, generar ansiedad en el usuario. La autonomía de un teléfono tiene datos objetivos, pero la experiencia de autonomía no equivale a la curva de carga de la batería. Lo que más angustia al usuario puede ser la espera de carga entre 95--100\%, o el estado de batería casi agotada entre 0--5\%; una autonomía brillante en el papel no significa una buena experiencia. Este caso muestra que los datos pueden decirte qué ocurrió, pero no te dicen automáticamente cómo lo percibe el usuario.

\begin{figure}[H]
\centering
\includegraphics[width=0.94\textwidth]{experience-not-data.png}
\caption{La experiencia no son los datos: los datos de autonomía y la experiencia de autonomía no son equivalentes. Diagrama conceptual propio, elaborado a partir de la conversación de 00:26:37--00:33:58.}
\end{figure}

\begin{knowledgebox}{Lectura de la figura: la experiencia combina datos y umbrales psicológicos}
A la izquierda, Data son indicadores como autonomía promedio, retención y clics; a la derecha, Experience son la ansiedad, los puntos sensibles y la percepción subjetiva. Al leer esta figura conviene notar que la experiencia no se opone a los datos, sino que exige saber qué intervalos de los datos son psicológicamente más sensibles para el usuario.
\end{knowledgebox}

\begin{knowledgebox}{Términos clave: datos, experiencia, retrovisor}
\noindent\begin{tabularx}{\textwidth}{p{0.18\textwidth}p{0.36\textwidth}X}
\toprule
Concepto & Significado & Juicio en este episodio \\
\midrule
Datos & Registro cuantificado del comportamiento pasado & Permiten verificar resultados, pero no generan una dirección por sí mismos. \\
Experiencia & Percepción subjetiva del usuario en umbrales concretos & Determina si el producto se entiende, genera confianza y se sigue usando. \\
Retrovisor & Los datos son como el espejo retrovisor al manejar & Ayudan a evaluar el trayecto recorrido, pero no pueden ver por ti el camino de enfrente. \\
Juicio de convicción & Los datos no lo respaldan a corto plazo, pero la dirección sigue mereciendo la apuesta & Es especialmente importante en la ventana de las aplicaciones de IA. \\
\bottomrule
\end{tabularx}
\end{knowledgebox}

\subsection{Las fortalezas de ByteDance y sus efectos secundarios}

Al pasar del entrenamiento en experiencia de OnePlus a ByteDance (字节), el problema se convierte en otro régimen de velocidad y escala. Esta subsección examina la segunda lección que le dio ByteDance: la toma de decisiones guiada por datos y la iteración de alta frecuencia. OnePlus lanza un teléfono al año; ByteDance puede publicar una versión cada dos semanas, y el gerente de producto debe moverse con rapidez entre datos, experimentos, crecimiento y coordinación organizacional. Ming Chaoping reconoce que ByteDance le permitió dominar métodos de producto a mayor escala, pero también señala los efectos secundarios: un sistema de datos muy fuerte puede sofocar las ideas creativas que surgen de un destello de inspiración, porque algo verdaderamente nuevo puede mostrar datos negativos a corto plazo.

\begin{warningbox}{La infraestructura de datos no sustituye el juicio creativo}
Si solo se persigue un resultado positivo en pruebas A/B de corto plazo, el equipo será cada vez más hábil para optimizar partes existentes, pero no necesariamente creará formas nuevas. Esto vale en especial para el emprendimiento de aplicaciones de IA: los usuarios nuevos, los grupos de personas nuevos y los flujos de trabajo nuevos rara vez tienen datos estables al principio.
\end{warningbox}

\subsection{La competencia de autos inteligentes: un entrenamiento temprano de reward model}

Antes de su carrera como gerente de producto, Ming Chaoping menciona que en sus dos últimos años de universidad prácticamente dormía en el laboratorio mientras preparaba la competencia de autos inteligentes. Esa experiencia se entiende bien en el contexto del emprendimiento de aplicaciones de IA: que el auto vaya más rápido, que esquive obstáculos con más estabilidad o que la comunicación sea más fluida se refleja de inmediato en el resultado de la competencia. Él lo llama un muy buen reward model: los cambios en el código, los sensores, el hardware y la colaboración del equipo regresan al resultado de una forma muy directa.

Esta experiencia explica por qué más adelante fue sensible al «entorno» y a la «retroalimentación». Muchos gerentes de producto solo toman decisiones frente a pantallas y tableros de datos, mientras que la competencia de autos inteligentes lo puso desde temprano en contacto con ciclos cerrados del mundo real: entrada de sensores, salida de control, bugs, velocidad y resultados de la competencia. Aunque YouWare es un producto de software, también busca un ciclo cerrado similar: el usuario expresa una intención, el sistema genera una obra interactiva, la obra se comparte y recibe retroalimentación, y en la siguiente ronda el producto se ajusta.

\begin{importantbox}{Experiencia práctica: una buena retroalimentación de producto debe parecerse al cronómetro de una competencia}
Si la retroalimentación de un sistema es demasiado lenta o difusa, al equipo le cuesta formar intuición. La retroalimentación intensa de la competencia de autos inteligentes genera adicción, y las aplicaciones de IA necesitan un mecanismo similar: si el usuario publica, si comparte, si otros interactúan con lo que hizo, si sigue creando; todo eso se acerca más a un reward real que un clic aislado.
\end{importantbox}

\subsection{Resumen de la sección}

La base de producto de Ming Chaoping proviene de cuatro entrenamientos: el debate entrenó la perspectiva de tercero, la competencia de autos inteligentes entrenó la retroalimentación en ciclo cerrado, OnePlus entrenó la sensibilidad a la experiencia y ByteDance entrenó los datos y la iteración. En conjunto, estas capacidades explican por qué más adelante insistió en el contenedor, el entorno y la per token valuation.
